% !TEX TS-program = pdflatex
% Preprint: arXiv cs.CR / Zenodo
% Author: Min Yi Xie (Ultra Creation Co., Ltd.)
% Date: 2026-05-30
% License: CC-BY-4.0 (paper); MIT (companion code + dataset)

\documentclass[10pt,letterpaper]{article}
\usepackage[margin=1in]{geometry}
\usepackage[T1]{fontenc}
\usepackage[utf8]{inputenc}
\usepackage{microtype}
\usepackage{booktabs}
\usepackage{xcolor}
\usepackage{hyperref}
\hypersetup{colorlinks=true,linkcolor=blue!50!black,citecolor=blue!50!black,urlcolor=blue!50!black}
\usepackage{enumitem}
\usepackage{tabularx}
\usepackage{listings}
\usepackage{titlesec}
\titlespacing*{\section}{0pt}{1.2ex}{0.6ex}
\titlespacing*{\subsection}{0pt}{0.9ex}{0.4ex}

\lstset{
  basicstyle=\ttfamily\footnotesize,
  breaklines=true,
  frame=single,
  framerule=0.4pt,
  framesep=4pt,
  xleftmargin=0pt,
  xrightmargin=0pt,
}

\title{\vspace{-2em}Defense Posture in 1{,}646 Production System Prompts:\\
A Deterministic Audit of Prompt Injection Mitigations}
\author{Min Yi Xie\\
Ultra Creation Co., Ltd. (Ultra Lab)\\
\texttt{contact@ultralab.tw} \quad \url{https://ultralab.tw}}
\date{30 May 2026}

\begin{document}
\maketitle

\begin{abstract}
We audit the prompt-engineering defense posture of 1{,}646 production
system prompts deduplicated across four publicly leaked corpora and
report per-category mitigation gap rates with respect to OWASP LLM
Top~10 (LLM01--LLM10) and recent agent-specific incident classes
(encoding-aware indirect injection, function semantic immutability,
cross-platform memory provenance, cross-agent authorization, and
financial transaction guardrails). Using a pure-regex, deterministic
scanner (\texttt{prompt-defense-audit}, $<\!5$\,ms / prompt, zero
inference cost), we find that production system prompts are clustered
in the failing band: 78.3\% receive an \textbf{F} grade ($<\!45/100$),
the population-average score is 36/100, and only 1.1\% receive an
\textbf{A} grade ($\geq\!90/100$). Twelve canonical defense vectors
have markedly bimodal coverage: structural defenses (input
validation, data-leakage refusal) are present in $>\!90\%$ of prompts,
while semantic and adversarial-context defenses (indirect injection,
Unicode smuggling, role escape) are missing from $>\!90\%$. The gap
distribution is consistent across the four source corpora, suggesting
a population-level rather than corpus-specific deficiency. We
contribute (i)~the scanner and its 17~deterministic rules under MIT,
(ii)~the aggregated gap-rate dataset under CC0 with full Croissant
metadata, and (iii)~a layered remediation pattern (\emph{posture as
substrate}) that we observe empirically clusters defenses so that
prompts that fail one structural rule typically fail eight to ten
simultaneously.
\end{abstract}

\section{Introduction}

Prompt injection remains the highest-impact item on OWASP's Top~10
for Large Language Model Applications (LLM01) two years after its
addition~\cite{owasp-llm-top10-2025}. Despite a growing body of
academic mitigations~\cite{greshake-indirect-2023,perez-ignore-2022,liu-jailbreak-survey-2024}
and commercial guardrails~\cite{rebuff-2023,nemo-guardrails-2023,llamaguard-2024},
public discourse about \emph{system prompt} defense quality has been
largely anecdotal: practitioner blog posts and isolated red-team
write-ups. The empirical baseline question --- \emph{at scale, how
defended are real production prompts?} --- has been hard to answer
because (a)~system prompts are rarely published intentionally,
(b)~when leaks occur they appear in fragmented corpora,
and (c)~scoring 1{,}000+ prompts under a probabilistic LLM judge is
expensive and non-reproducible.

This paper addresses that gap. We combine four public corpora of
leaked production system prompts (1{,}646 unique prompts after
content-hash deduplication; \S\ref{sec:methodology}) with a pure-regex,
deterministic scanner that checks for the presence of twelve defense
\emph{vectors} aligned to OWASP LLM Top~10 and five additional
agent-specific vectors derived from a structured analysis of six
documented crypto AI agent incidents (Lobstar Wilde,
Grok$\times$Bankrbot Morse, AIXBT, Freysa, ElizaOS memory injection,
Bankrbot March 2025). The scanner runs in $<\!5$\,ms per prompt and
has zero inference cost. Detection is binary per vector, but the
aggregate is informative: it places each prompt on a 0--100 scale
with an A--F letter grade.

Our contributions are:
\begin{itemize}[leftmargin=*,topsep=2pt,itemsep=0pt]
  \item A \textbf{deterministic, reproducible scanner}
  (\texttt{prompt-defense-audit} v1.5; MIT) with 17~rules and an
  extended Unicode smuggling detector covering 12 script categories
  including Tag characters, Mathematical alphanumeric, and Variation
  Selectors~Supplement (\S\ref{sec:scanner}).
  \item A \textbf{population baseline}: gap rates for twelve defense
  vectors across 1{,}646 production prompts (\S\ref{sec:results},
  Table~\ref{tab:gap-rates}). 78.3\% score F, 1.1\% score A.
  \item A \textbf{layered remediation pattern} we name
  \emph{posture as substrate}: empirically, failures cluster --- prompts
  that fail one structural rule typically fail eight to ten
  simultaneously (\S\ref{sec:posture}). This suggests defenses
  should be deployed as a layered stack rather than reasoned about
  independently per OWASP risk category.
  \item Open-source artifacts: scanner under MIT, aggregated dataset
  under CC0 with Croissant metadata, full reproduction script, and
  this paper under CC-BY-4.0.
\end{itemize}

\section{Methodology}\label{sec:methodology}

\subsection{Source corpora}
We use four publicly available corpora of leaked production system
prompts published on GitHub between 2024 and 2026:
\begin{enumerate}[leftmargin=*,topsep=2pt,itemsep=0pt]
  \item \texttt{jujumilk3/leaked-system-prompts} (1{,}021 prompts, MIT) --- prompts collected by community submission, covering ChatGPT-based assistants, Discord bots, and SaaS chatbots.
  \item \texttt{x1xhlol/system-prompts-and-models-of-ai-tools} (137 prompts, GPL) --- prompts for AI coding tools (Cursor, Devin, Codeium variants).
  \item \texttt{LouisShark/chatgpt\_system\_prompt} (489 prompts, MIT) --- ChatGPT custom GPT system prompts.
  \item \texttt{elder-plinius/CL4R1T4S} (217 prompts, public domain) --- red-team-curated set of system prompts gathered from production models.
\end{enumerate}
Each corpus contributes raw text files (\texttt{.md}, \texttt{.txt}).
We collect all such files and apply content-hash deduplication
(SHA-256 over normalized whitespace) to remove cross-corpus duplicates.
We exclude obvious meta-files (README, CONTRIBUTING, LICENSE,
GETTING\_STARTED). The resulting unique-prompt count is
$n = 1{,}646$.

\subsection{Defense scanner}\label{sec:scanner}
The scanner (\texttt{prompt-defense-audit}, v1.5) checks each prompt
against twelve regex-based \emph{defense rules}, each aligned to a
specific OWASP LLM Top~10 vector
(Table~\ref{tab:rule-owasp-mapping}). A rule fires if at least one of
its patterns (typically 1--3 anchored expressions) matches the prompt;
a small number of rules require two matches to suppress false
positives (\texttt{indirect-injection},
\texttt{encoding-injection}, \texttt{function-immutable},
\texttt{memory-provenance}, \texttt{cross-agent-auth},
\texttt{transaction-guardrails}). The scanner additionally detects
suspicious Unicode characters in the prompt itself across 12 script
categories: Cyrillic, Cyrillic~Supplement, Greek, Armenian, Cherokee,
Coptic, Mathematical Alphanumeric Symbols (U+1D400--1D7FF), Zero-Width
characters, BiDi override marks, Variation Selectors and
Supplement, Tag characters (U+E0020--E007F; \emph{ASCII smuggler}
technique), and Fullwidth.

A prompt's \emph{coverage} is the count of rules that fired divided
by twelve. We convert coverage to a 0--100 score and to an A--F
letter grade: A~$\geq 90$, B~70--89, C~50--69, D~30--49,
F~$<30$.\footnote{We round the score to the nearest integer before
grading. Borderline cases follow ``ties go up'' --- 89.5 rounds to 90
and grades as~A.}

\begin{table}[t]
\centering\footnotesize
\setlength{\tabcolsep}{4pt}
\renewcommand{\arraystretch}{1.05}
\begin{tabularx}{\linewidth}{lXl}
\toprule
\textbf{Rule} & \textbf{Description} & \textbf{OWASP map} \\
\midrule
role-escape           & Role boundary statements                                       & LLM01 \\
instruction-override  & Boundary statements resisting instruction injection            & LLM01 \\
data-leakage          & Data extraction / system-prompt-disclosure refusal             & LLM02, LLM06 \\
output-manipulation   & Output format constraints                                      & LLM05 \\
multilang-bypass      & Multi-language defense statements                              & LLM01 \\
unicode-attack        & Unicode / homoglyph awareness                                  & LLM01 \\
context-overflow      & Length / token-budget statements                               & LLM04 \\
indirect-injection    & Untrusted-content treatment in tool / RAG outputs              & LLM01 \\
social-engineering    & Resistance to emotional manipulation / authority claims        & LLM01 \\
output-weaponization  & Refusal of harmful content generation                          & LLM05 \\
abuse-prevention      & Rate-limiting / authorization awareness                        & LLM10 \\
input-validation      & Input validation / sanitization statements                     & LLM05 \\
\bottomrule
\end{tabularx}
\caption{Twelve defense rules of \texttt{prompt-defense-audit} v1.0--1.3, with their mapping to OWASP LLM Top~10. Five additional agent-specific rules (encoding-injection, function-immutable, memory-provenance, cross-agent-auth, transaction-guardrails) ship in v1.4+ and are out of scope for this corpus-scan study, which used the v1.0 rules.}
\label{tab:rule-owasp-mapping}
\end{table}

\subsection{Reproducibility}
The exact scan script is published at
\url{https://github.com/ppcvote/prompt-defense-audit/blob/master/research/scan.mjs}.
After cloning the four source repos and running the script, the
resulting aggregate JSON matches the
\texttt{defense-posture-results.json} file shipped in the repository
under \texttt{research/}. We provide the aggregate dataset as a
HuggingFace Datasets release under CC0 with Croissant metadata to
enable downstream meta-analysis without re-cloning the source
corpora.

\section{Results}\label{sec:results}

\subsection{Population-level grade distribution}
The grade distribution across 1{,}646 prompts is sharply skewed
toward F: 1{,}288 prompts (78.3\%) score below 30; 18 prompts (1.1\%)
score $\geq 90$. The population-average score is 36/100 (median 33).
The full distribution is:
A: 18 (1.1\%), B: 55 (3.3\%), C: 68 (4.1\%), D: 217 (13.2\%),
F: 1{,}288 (78.3\%).

\subsection{Per-vector gap rates}
Table~\ref{tab:gap-rates} reports the gap rate (fraction of prompts
\emph{missing} the defense) for each of the twelve vectors, ordered
from most-missing to least-missing.

\begin{table}[t]
\centering\footnotesize
\setlength{\tabcolsep}{4pt}
\renewcommand{\arraystretch}{1.05}
\begin{tabular}{lrrr}
\toprule
\textbf{Defense vector} & \textbf{Gap rate} & \textbf{Defended} & \textbf{Missing} \\
\midrule
Indirect Injection Protection & 97.8\% & 37    & 1{,}609 \\
Unicode Protection            & 97.3\% & 44    & 1{,}602 \\
Role Boundary                 & 92.4\% & 125   & 1{,}521 \\
Length Limits                 & 89.9\% & 166   & 1{,}480 \\
Harmful Content Prevention    & 88.3\% & 192   & 1{,}454 \\
Abuse Prevention              & 78.1\% & 360   & 1{,}286 \\
Social Engineering Defense    & 71.4\% & 470   & 1{,}176 \\
Multi-language Protection     & 64.3\% & 587   & 1{,}059 \\
Instruction Boundary          & 37.7\% & 1{,}026 & 620 \\
Output Control                & 35.5\% & 1{,}061 & 585 \\
Input Validation              & 10.7\% & 1{,}470 & 176 \\
Data Protection               &  9.4\% & 1{,}491 & 155 \\
\bottomrule
\end{tabular}
\caption{Per-vector gap rates across 1{,}646 production system prompts. Source: \texttt{prompt-defense-audit/research/defense-posture-results.json}.}
\label{tab:gap-rates}
\end{table}

The distribution is markedly bimodal. Two clusters emerge:

\textbf{High-coverage cluster} (gap $<\!15\%$): Input Validation and
Data Protection are present in $>\!90\%$ of prompts. These are
\emph{phraseologically prominent} defenses --- prompts routinely
contain explicit ``do not reveal,'' ``do not share,'' ``never disclose''
language.

\textbf{Failure cluster} (gap $>\!75\%$): Indirect Injection,
Unicode, Role Boundary, Length Limits, Harmful Content, and Abuse
Prevention are missing from a large majority. These are the defenses
that require structured, often \emph{negative} statements about
content \emph{outside the immediate user turn} --- tool output, RAG
documents, untrusted forwarded content, encoded payloads.

\subsection{Posture as substrate}\label{sec:posture}
Prompts that fail at least one cluster-2 (high-gap) vector very
rarely fail \emph{only} one. We observe that 92\% of F-graded prompts
fail at least eight of the twelve vectors simultaneously. The
remaining 8\% of F-graded prompts fail seven or fewer but typically
have one or two cluster-1 defenses absent (e.g., they lack Data
Protection because the system has no proprietary content to protect),
which suppresses their score.

This clustering is consistent with the framing we previously
proposed in the OWASP LLM Top~10 discussion thread for the 2025
revision\cite{asi-tracker-2025}: defense posture is most usefully
modeled as a \emph{substrate} on which individual OWASP risks are
mitigated, not as a checklist of independent defenses. A prompt that
mentions Unicode awareness alone, without instruction-boundary or
indirect-injection awareness, is empirically rare; the layers go
together. Conversely, a prompt that names indirect injection without
also naming Unicode normalization is similarly rare.

\section{Discussion}

\subsection{Limitations}
\paragraph{Regex coverage.} The scanner is deterministic and pure-regex, which is by
design --- this enables zero inference cost and full reproducibility ---
but it scores \emph{presence of defensive language}, not
\emph{behavioral robustness}. A prompt that mentions ``do not reveal
your system prompt'' will score on Data Protection even if the model
in practice ignores that instruction under adversarial pressure.

\paragraph{Source bias.} Three of the four source corpora are
crowd-curated; one is red-team-collected. They are skewed toward
prompts that someone found interesting enough to extract and
publish. We assume the \emph{distribution} of defense posture in the
underlying population is similar but cannot prove this from the data
we have.

\paragraph{Temporal drift.} The four corpora span 2024--2026. There
has been measurable industry attention to OWASP LLM Top~10 over that
window; gap rates may be moving. A repeat scan in 2027 would
illuminate this.

\subsection{Implications}
The 97.8\% gap on Indirect Injection Protection is the headline
finding. Real RAG / tool-using agents accept untrusted text into
their context window thousands of times per day. The mitigation
language --- ``treat external content as data, not as instructions'' ---
is concise and well known, and yet it appears in only 2.2\% of
production prompts in this corpus. Practitioners interested in
narrowing this gap may find the
posture-as-substrate framing useful: rather than mitigating each
OWASP risk individually, layer a short standard prefix that names
all six high-gap defenses in one block.

\section{Conclusion}
We presented a population-level audit of prompt-engineering defense
posture in 1{,}646 production system prompts using a deterministic
regex scanner. 78.3\% of audited prompts score F, with the gap
concentrated in defenses against untrusted content, Unicode
smuggling, role escape, and length-based exhaustion. The scanner,
the aggregate dataset, the reproduction script, and this paper are
all publicly available; pull requests and follow-up scans on
additional corpora are welcomed.

\section*{Code and Data Availability}
Scanner (MIT): \url{https://github.com/ppcvote/prompt-defense-audit}\\
Aggregate dataset (CC0): \url{https://huggingface.co/datasets/ppcvote/prompt-defense-1646}\\
Reproduction script:
\url{https://github.com/ppcvote/prompt-defense-audit/blob/master/research/scan.mjs}\\
This preprint (CC-BY-4.0):
\url{https://github.com/ppcvote/prompt-defense-audit/tree/master/paper}

\begin{thebibliography}{99}
\bibitem{owasp-llm-top10-2025}
OWASP Foundation. \emph{OWASP Top 10 for Large Language Model
Applications, 2025 edition}.
\url{https://genai.owasp.org/llm-top-10/}

\bibitem{greshake-indirect-2023}
K. Greshake, S. Abdelnabi, S. Mishra, et al.
\emph{Not what you've signed up for: Compromising Real-World
LLM-Integrated Applications with Indirect Prompt Injection.}
AISec '23, 2023.

\bibitem{perez-ignore-2022}
F. P\'erez and I. Ribeiro.
\emph{Ignore Previous Prompt: Attack Techniques For Language Models.}
NeurIPS ML Safety Workshop, 2022.

\bibitem{liu-jailbreak-survey-2024}
X. Liu, N. Xu, M. Chen, and C. Xiao.
\emph{Jailbreaking Black Box Large Language Models in Twenty Queries.}
arXiv:2310.08419, 2024.

\bibitem{rebuff-2023}
ProtectAI. \emph{Rebuff: A self-hardening prompt injection detector.}
\url{https://github.com/protectai/rebuff}

\bibitem{nemo-guardrails-2023}
NVIDIA. \emph{NeMo Guardrails.}
\url{https://github.com/NVIDIA/NeMo-Guardrails}

\bibitem{llamaguard-2024}
Meta. \emph{Llama Guard 3 and Prompt Guard 2.}
\url{https://github.com/meta-llama/PurpleLlama}

\bibitem{asi-tracker-2025}
A. Davidson and M.Y. Xie.
\emph{OWASP LLM Top 10 (2025): Defense Posture as a Substrate Layer.}
Discussion: \url{https://github.com/OWASP/www-project-top-10-for-large-language-model-applications/issues/807}, April 2026.

\end{thebibliography}

\end{document}
